%% ECON 282E -- Homework 2, Part B: Heterogeneity and Learning (S8-S9, route, audit)
%% RE-CHECK WHEN THE S8-S9 DECKS ARE BUILT.  This part was drafted on 2026-09-26
%% against the AI-ECON-2026 source decks that S8-S9 will draw on -- Session03 (RL:
%% TD error, deadly triad, actor-critic), Session04 (HA: Krusell-Smith, OLG
%% transitions, DEQN/DeepHAM) and Session05 (text: the validation ladder,
%% non-classical measurement error, generated regressors).  Deck numbers cited
%% below (8.A, 8.B, 9.A, 9.B) are block-level only until the decks exist.
%% Build:  python3 tools/build_docs.py homeworks-exams/HW2/hw2b.tex
%% With answers:  python3 tools/build_docs.py --solutions homeworks-exams/HW2/hw2b.tex
%% Answers live in solutions/ (gitignored); solution PDFs are gitignored too.

\documentclass[11pt]{article}

\newif\ifsolutions \solutionsfalse

\usepackage[margin=1in]{geometry}
\usepackage{amsmath,amssymb,booktabs,enumitem,xcolor}
\usepackage[T1]{fontenc}
\usepackage{microtype}
\usepackage[colorlinks=true,linkcolor=blue!50!black,urlcolor=blue!50!black]{hyperref}

\newcommand{\solution}[1]{\ifsolutions\medskip
  {\color{blue!45!black}\textbf{Solution.} #1}\medskip\fi}

\AtBeginDocument{\ifsolutions
  \IfFileExists{solutions/hw2b_key.tex}{\input{solutions/hw2b_key.tex}}
    {\GenericWarning{}{Answer key solutions/hw2b_key.tex not found}}\fi}

\title{\vspace{-2em}\textbf{ECON 282E --- Homework 2, Part B}\\[2mm]
       \large Heterogeneity and Learning: Two New Tools, Your Route, and an Audit}
\author{Foundations of Macroeconomics \\ University of California, Riverside}
\date{Out: Thursday, November 12, 2026 \quad\textbullet\quad Due (with Part A): Monday, November 23, 2026}

\begin{document}
\maketitle
\thispagestyle{empty}

\noindent\rule{\linewidth}{0.4pt}\par
\small
\textbf{What Part B adds.} Two short core items put the last two families of tools in your hands ---
reinforcement learning (Session 8) and text as data (Session 9). Then the \textbf{route}: one of four
directions, in which \emph{you} pose the question and choose the combination of models and methods.
Finally, the \textbf{audit}: you check a classmate's HW1 the way a referee or a co-author would.

\textbf{Timing.} C3 uses Session 8 (November 12); C4 and route H4 use Session 9 (November 19). Plan
your route in the charter now; the peer audit can be done as soon as it is assigned.

\textbf{Your final project.} Your team (formed November 5; proposal November 16) will assemble at
least one HW1 component and one HW2 component. Choose your route with that in mind, and say in the
charter which piece of the project it is meant to become.
\normalsize\par
\noindent\rule{\linewidth}{0.4pt}\par

%=============================================================================
\section*{C3. Reinforcement learning, minimal (8 points)}
\emph{Session 8, Block A: dynamic programming without the model.}

Take the $\delta=1$ growth model of HW1 A2 on a 40-point capital grid, where the choice $k'$ is also a
grid point. VFI on this grid is your benchmark, and $g^*=\alpha\beta k^\alpha$ your truth.
\begin{enumerate}[label=(\alph*), leftmargin=*]
  \item \textbf{(5 pts)} Implement tabular Q-learning: the agent never sees the transition or reward
        functions, only sampled $(s,a,r,s')$. State your exploration rule and learning-rate schedule.
        For sample budgets $10^4,10^5,10^6$ (and more if you can), report the share of states where
        the greedy policy equals VFI's, $\max|g/g^*-1|$, $\max_s|\max_aQ(s,a)-V(s)|$, and how many
        state--action pairs were never visited.
  \item \textbf{(3 pts)} Count the cost: how many \emph{samples} did Q-learning need, against how many
        \emph{Bellman evaluations} VFI needed, to reach the same policy? When is that trade worth
        making in economics --- what does RL buy that VFI cannot?
\end{enumerate}

\solution{\keyDone}

%=============================================================================
\section*{C4. Text as measurement, minimal (8 points)}
\emph{Session 9.}

The folder \texttt{homeworks-exams/HW2/data/fomc/} holds every FOMC post-meeting statement,
1994--2026 (public domain; loader \texttt{fomc\_data.py}).
\begin{enumerate}[label=(\alph*), leftmargin=*]
  \item \textbf{(4 pts) Build an index.} Construct one index per statement --- an uncertainty index,
        a hawk--dove stance score, or one of your own design --- by a dictionary, embeddings, or a
        language model. Document every choice so that someone else gets the same number.
  \item \textbf{(4 pts) Climb the validation ladder.} Validate it against a public series you download
        yourself (e.g.\ \texttt{VIXCLS} or the change in \texttt{DFF} after the meeting, from FRED):
        correlation, a placebo (the index shuffled in time), stability across sub-samples, and five
        statements read by hand. Then say how the index could discipline one parameter of your C1
        economy, and why its measurement error is \emph{not} the classical kind (9.A).
\end{enumerate}

\solution{\keyDtwo}

%=============================================================================
\section*{Your route (24 points)}

Choose \textbf{one}. Each route asks \emph{you} to pose a question, and each one grades the same
four things: \textbf{(6)} the question, stated as a project manager would (positive or normative,
and what number answers it); \textbf{(8)} the model and method, and why this combination; \textbf{(6)}
validation against something you already trust --- your C1 solution, your HW1 engine, or a
published number; \textbf{(4)} what you would do next, and which piece of your final project this
becomes.

\begin{description}[leftmargin=0pt, itemsep=4pt]
\item[H1. Aggregate risk and who bears it] \emph{(Sessions 6, 8.)} Add aggregate TFP risk to C1 and
  solve Krusell--Smith: a perceived law of motion $H$, simulate, regress, update. Report the
  forecast rule, its $R^2$ \emph{and} the Den Haan statistic. Then return to HW1's normative question
  with heterogeneity: who --- by wealth or earnings --- bears the cost of business cycles?
  Optionally check one object with a deep method (DEQN or DeepHAM, 8.B).
\item[H2. Fiscal policy and the transition] \emph{(Sessions 2, 6, 8.)} A tax or debt reform in your
  C1 economy or in an OLG economy: compare the two steady states, compute the transition between
  them by shooting on the capital path, and report welfare in consumption-equivalent terms by wealth
  decile (or by cohort). Does the steady-state comparison get the sign right for everyone?
\item[H3. A learning method, pushed] \emph{(Sessions 7, 8.)} Take one AI solver further than C2:
  deep VFI or actor--critic on the C1 household, RL on a problem where the transition is unknown to
  the agent, or a network that takes the distribution as input. Validate against C1. Name its
  failure modes, and demonstrate at least one (for example the deadly triad, or a residual that is
  small where the data are and large where the policy matters).
\item[H4. Text-disciplined heterogeneity] \emph{(Sessions 1, 9.)} Measure a shock from text or other
  public data --- AI exposure of occupations, an uncertainty index, a policy stance --- and feed it into
  the C1 economy, for instance as a change in the earnings process. What does the measurement imply for
  inequality or for $r^*$, and how much of that answer is measurement error?
\end{description}

\solution{\keyDthree}

%=============================================================================
\section*{Peer replication audit (10 points)}

You will be sent, privately, a classmate's anonymized HW1 package. Audit it as a referee would, in at
most one page (template \texttt{audit.md} beside this sheet):
\begin{enumerate}[label=(\alph*), leftmargin=*]
  \item Does \texttt{run\_all.py} run on a clean environment? What did you have to fix?
  \item Reproduce \textbf{three} numbers from the paper. Which match, to how many digits?
  \item Is the validation convincing? Name the single check you would ask the author to add.
  \item One thing the author did better than you did.
\end{enumerate}
The audit is graded for \emph{you}, the reviewer: on accuracy, specificity and tone. The author
receives it as feedback only; it does not change their HW1 grade. Do not identify the author, or
yourself, in the report.

\solution{\keyDfour}

\vfill
\noindent\rule{\linewidth}{0.4pt}\par
\small
\textbf{Hand in, with Part A:} the PDF, the repository with \texttt{run\_all.py}, the updated charter,
\texttt{AI\_LOG.md} with its validation ledger, and the audit. \textbf{Part B: 50 points.}
\normalsize

\end{document}
